\documentclass[11pt]{article}

\usepackage{amsmath,amssymb}
\usepackage{graphicx}
\usepackage{hyperref}
\usepackage{geometry}
\usepackage{tikz}

\geometry{margin=2.5cm}

\title{
	A Dirac Operator from the Logarithmic Geometry of Primes
}

\author{
	Thomas Hoffbauer
}

\date{\today}

\begin{document}
	
	\maketitle
	
	\begin{abstract}
		We propose a Dirac-type operator constructed from the logarithmic
		geometry of prime numbers.
		Numerical experiments suggest that the spectrum of this operator
		is strongly correlated with the ordinates of the nontrivial zeros
		of the Riemann zeta function.
		This observation motivates a spectral conjecture
		in the spirit of the Hilbert--Pólya program.
	\end{abstract}
	
	\section{Introduction}
	
	The nontrivial zeros of the Riemann zeta function
	\[
	\zeta(s)
	\]
	play a fundamental role in number theory.
	
	The famous Riemann Hypothesis states that all nontrivial zeros lie on
	the critical line
	\[
	\mathrm{Re}(s)=\tfrac12 .
	\]
	
	A long-standing idea, often attributed to Hilbert and Pólya,
	is that the imaginary parts of the zeros might appear
	as eigenvalues of a self-adjoint operator.
	
	The aim of this paper is to suggest a simple candidate operator
	constructed directly from the geometry of prime numbers.
	
	\section{Logarithmic Geometry of Primes}
	
	Let
	\[
	p_1,p_2,p_3,\dots
	\]
	be the sequence of prime numbers.
	
	We introduce the logarithmic coordinate
	
	\[
	C = \log p.
	\]
	
	This transformation has two important properties:
	
	\begin{itemize}
		\item multiplicative relations become additive,
		\item growth patterns become approximately linear.
	\end{itemize}
	
	The primes therefore form a discrete set in a logarithmic space.
	
	\section{Prime Graph}
	
	We define a weighted graph
	\[
	G=(V,E)
	\]
	with
	
	\[
	V=\{p_1,p_2,\dots,p_N\}.
	\]
	
	Between two primes we introduce the symmetric weight
	
	\[
	w(p_i,p_j)
	=
	\frac{\log(p_i p_j)}{\sqrt{p_i p_j}}.
	\]
	
	This coupling combines
	
	\begin{itemize}
		\item logarithmic geometry,
		\item scale damping.
	\end{itemize}
	
	\section{Dirac-Type Operator}
	
	Let
	
	\[
	Q_{ij}=w(p_i,p_j).
	\]
	
	We define the operator
	
	\[
	D=
	\begin{pmatrix}
		0 & Q\\
		Q^T & 0
	\end{pmatrix}.
	\]
	
	This matrix is self-adjoint and therefore has a real spectrum
	
	\[
	D\psi_n=\lambda_n\psi_n .
	\]
	
	\section{Numerical Experiments}
	
	Using the first primes and the first Riemann zeros
	\[
	\frac12+i\gamma_n
	\]
	we compare the sequences
	
	\[
	\lambda_n
	\quad \text{and} \quad
	\gamma_n .
	\]
	
	Numerical experiments indicate a strong monotonic correlation.
	
	For moderate matrix sizes the correlation coefficient
	is approximately
	
	\[
	\mathrm{corr} \approx 0.97 .
	\]
	
	This suggests a possible spectral relationship.
	
	\section{Spectral Conjecture}
	
	We propose the following conjecture.
	
	\textbf{Bamberg Spectral Conjecture}
	
	There exists a self-adjoint operator
	\[
	D_{BM}
	\]
	constructed from the logarithmic geometry of primes
	such that
	
	\[
	\mathrm{Spec}(D_{BM})
	=
	\{\gamma_n\},
	\]
	
	where $\gamma_n$ are the ordinates of the nontrivial
	zeros of the Riemann zeta function.
	
	\section{Interpretation}
	
	The proposed construction connects several structures:
	
	\begin{center}
		\begin{tabular}{c|c}
			structure & interpretation \\
			\hline
			primes & discrete geometry \\
			operator & dynamics \\
			eigenvalues & spectrum \\
			zeta zeros & resonances
		\end{tabular}
	\end{center}
	
	In this picture the zeros of the zeta function
	appear as spectral quantities of an operator
	determined by the primes.
	
	\section{Outlook}
	
	Several questions remain open:
	
	\begin{itemize}
		\item the precise operator form,
		\item asymptotic spectral statistics,
		\item connections to trace formulas.
	\end{itemize}
	
	Further numerical investigations with larger matrices
	may provide additional insight.
	
	\section*{References}
	
	Riemann, B. (1859).  
	On the number of primes less than a given magnitude.
	
	Montgomery, H.  
	The pair correlation of zeros of the zeta function.
	
	Odlyzko, A.  
	Tables of Riemann zeta zeros.
	
\end{document}